\section{Construction}\label{sec:construction}

The index is built from one all-size peel and one pass per size, without ever holding all core values at once.

\stitle{Core values.}
We reuse the all-size vertex peel of the clique-tree line \cite{Pivoter}: a pivoting clique tree is built once, every vertex's number of $s$-cliques is a sum of binomial coefficients over the tree's leaves, and a bucket peel per size removes the vertex of smallest count and updates the counts of the affected leaves in closed form.
Sizes share work: the previous size's order and values bound the next size's, and a vertex whose value has reached its floor (Theorem~\ref{thm:tail}) is settled for all larger sizes.
The peel delivers the values of one size at a time; the index keeps a two-row window, not the $s_{\max}\times n$ matrix.

\stitle{One tree per size, on arrival.}
When the values of size $s$ arrive, the vertices are activated from the highest value down.
A leaf of the clique tree becomes live when its mandatory vertices are active and at least $s$ of its vertices are active; the active vertices of a live leaf are pairwise adjacent inside a common clique, so they are joined in a union-find structure.
After each level the components that changed are the canonical nodes created at that level, and every activated vertex records its own node.
This is the pass that produces $\Tree{s}$ with tops, parents and own nodes.

\stitle{Chains by refinement.}
Chains are refined size by size in a prefix trie: at $s=2$ vertices are grouped by their own node; at each later size every group is split by the own node of that size; a vertex whose $\omega$ has been reached terminates at its current trie node.
The trie's DFS order, terminating chains before children, is the lexicographic key order of Section~\ref{sec:order}, so it yields the chain ranks, the aligned labels and the bitmap directly.
Working memory beyond the peel is one trie node per distinct own-node prefix and $O(n)$ words.

\stitle{Run arrays and values.}
For each size, the tree is traversed with children ordered by their smallest chain rank; the resulting chain sequence is merged into runs and every node's entry is recorded.
Each chain then receives $\omega$, $\sigma$, its residue and its trajectory, and the per-size widths are chosen from the maxima.

\stitle{Cost.}
The peel and the clique tree dominate; the trie, the ranks, the run arrays and the values add a few percent (Section~\ref{sec:exp}).
Peak memory is the peel's clique tree plus the graph; the chain index itself is small.
